% B2: Cross-Domain Transfer and Feature Pipeline
% Documents why PCA trained on RouteLLM generalizes to LMSYS, and the
% complete feature extraction pipeline used by the production router.
% Cross-references: Figure 1 (main paper), B1 (validation methodology),
%                   Table 2 (data provenance), F1 (configuration details)

\subsection{Cross-Domain Transfer and Feature Pipeline}
\label{sec:cross_domain_transfer}

A key methodological choice in banditGPT is training the PCA on one dataset
(RouteLLM battles) and evaluating on another (LMSYS Chatbot Arena prompts).
This subsection explains why this cross-domain transfer works, documents the
complete feature pipeline, and discusses the limitations of the current
approach.

\subsubsection{Data Provenance and Independence}

\paragraph{Training and evaluation datasets.}
Table~\ref{tab:data_provenance} summarizes the two data sources.

\begin{table}[h]
\centering
\caption{Data provenance: training and evaluation are independent by design.}
\label{tab:data_provenance}
\small
\begin{tabular}{lll}
\toprule
 & \textbf{PCA Training (Warmup)} & \textbf{Evaluation (Holdout)} \\
\midrule
\textbf{Source} & RouteLLM HuggingFace dataset & LMSYS Chatbot Arena \\
\textbf{Size} & 80,000 prompts & 750 prompts \\
\textbf{Collection} & Curated battle pairs & General user submissions \\
\textbf{Time Period} & Pre-2024 collection & Independent collection \\
\textbf{Model Pair} & Mixtral vs.\ GPT-4-Turbo & Mixtral vs.\ GPT-4-Turbo \\
\textbf{Overlap} & \multicolumn{2}{c}{243 incidental duplicates removed (0.24\%)} \\
\bottomrule
\end{tabular}
\end{table}

The datasets are disjoint by provenance: they were collected independently by
different projects with different sampling procedures.  The 243 incidentally
overlapping prompts (detected via exact string matching, representing 0.24\% of
the warmup set) were removed to eliminate any residual contamination.

\paragraph{Why the same model pair matters.}
Cross-domain transfer succeeds because both datasets involve the same
underlying models.  The PCA captures variance in how prompt embeddings relate
to the Mixtral-vs-GPT-4-Turbo capability boundary.  This capability boundary
is a property of the \emph{models}, not the prompt distribution---a
``write me a poem'' prompt activates the same model strengths regardless of
whether it was collected by RouteLLM or LMSYS.

A different model pair (e.g., Claude-3-Haiku vs.\ GPT-4o) would define a
different capability boundary and would likely require retraining the PCA.
This is a limitation of the current two-model evaluation topology
(Section~\ref{sec:transfer_limitations}).

\subsubsection{Feature Extraction Pipeline}

The production router extracts features through the following pipeline,
implemented in \texttt{FeatureService} (\texttt{src/bandit\_gpt/feature\_service.py}):

\paragraph{Step 1: Sentence embedding.}
Each prompt is encoded by a sentence transformer
(\texttt{all-MiniLM-L6-v2}~\cite{reimers2019sentencebert}) into a
384-dimensional L2-normalized embedding vector.  This model was chosen for its
balance of quality and inference speed (${\sim}5$ms per prompt on CPU).

\paragraph{Step 2: PCA compression.}
The 384-dimensional embedding is projected to 32 dimensions via PCA trained on
the 80K RouteLLM warmup prompts.  The 32 components capture 35.14\% of the
total embedding variance.  While this may appear low in absolute terms,
the remaining ${\sim}$65\% represents variation orthogonal to model preference
(e.g., topic, language, style) that is irrelevant for routing.
Figure~1 validates that the variance captured by PC1 correlates
significantly with the reward gap ($p < 0.0001$), confirming that the
PCA concentrates routing-relevant signal.

\paragraph{Step 3: Bias term.}
A constant bias term (1.0) is appended, yielding a 33-dimensional feature
vector: $[PCA_0, PCA_1, \ldots, PCA_{31}, \text{bias}]$.  The bias allows
LinUCB to learn a per-model intercept independent of prompt content.

\paragraph{Summary.}
The complete pipeline is:
\begin{equation}
\text{prompt} \xrightarrow{\text{MiniLM}} \mathbb{R}^{384}
\xrightarrow{\text{PCA}} \mathbb{R}^{32}
\xrightarrow{[\cdot\,;\, 1]} \mathbb{R}^{33}
\label{eq:feature_pipeline}
\end{equation}

\subsubsection{PCA Variance and Routing Signal}

\paragraph{Explained variance breakdown.}
Table~\ref{tab:pca_variance} shows the cumulative explained variance for
selected component counts.

\begin{table}[h]
\centering
\caption{PCA explained variance (trained on 80K RouteLLM prompts).}
\label{tab:pca_variance}
\small
\begin{tabular}{lcc}
\toprule
\textbf{Components} & \textbf{Cumulative Variance} & \textbf{Feature Dim} \\
\midrule
1 (PC1 only) & 3.10\% & 2 \\
2 (PC1+PC2) & 5.40\% & 3 \\
23 & 29.01\% & 24 \\
\textbf{32 (production)} & \textbf{35.14\%} & \textbf{33} \\
384 (full) & 100\% & 385 \\
\bottomrule
\end{tabular}
\end{table}

\paragraph{Why 35\% variance is sufficient.}
The PCA operates on general-purpose sentence embeddings that encode
\emph{all} semantic content (topic, intent, style, language, length).
Only a fraction of this variance is relevant for routing---the component
that distinguishes ``prompts where GPT-4-Turbo wins'' from ``prompts where
Mixtral wins.''  The key validation is not explained variance \emph{per se},
but whether the captured components predict model preference.  Figure~1
confirms this: PC1 alone (3.10\% variance) achieves $p < 0.0001$ correlation
with the reward gap, and the full 32-component vector provides at least as
much signal.

\paragraph{Why not more components?}
Additional components beyond 32 add diminishing routing signal while
increasing LinUCB's per-update cost ($O(d^2)$ for Sherman-Morrison updates).
The 32-component choice balances signal capture against computational
efficiency.  With $d = 33$ (including bias), each LinUCB update requires
${\sim}$1,089 multiply-accumulate operations---fast enough for real-time
routing at 100+ QPS.

\subsubsection{Multi-Model Evaluation Data ($K{>}2$)}
\label{sec:multimodel_data}

The two-model experiments (Figures~1--5) use PCA trained on 80K RouteLLM
battles.  For the multi-model experiments (Figures~6, 8, and~9;
$K{=}5$ and $K{=}10$), we extend the data pipeline as follows.

\paragraph{43-model evaluation dataset.}
Rewards for 43 models are collected on the same LMSYS prompt pool using
GPT-4o as pairwise judge (Section~\ref{appendix:reward_signal}).  Only
prompts with full 43-model coverage are retained.

\paragraph{Three-way split.}
To avoid data leakage, the dev prompts with 43-model coverage are split into
two disjoint subsets via stratified sampling (seed~42):
\begin{itemize}[leftmargin=*,noitemsep,topsep=0pt]
    \item \textbf{Prior-training set} (${\sim}40\%$): Used to build LinUCB
          warmup priors for all 43 models.  Each (prompt, model) pair
          contributes a rank-1 update ($A_m \mathrel{+}= x x^\top$,
          $b_m \mathrel{+}= r \cdot x$), then matrices are scaled by a
          plasticity factor of~0.1.
    \item \textbf{Online-learning set} (${\sim}60\%$, 533 prompts): Fed to the
          bandit during evaluation as streaming interactions.
\end{itemize}
The holdout set (750~prompts) is never touched during prior construction or
online learning.  Pairwise disjointness of all three splits is verified
programmatically before each experiment.

\paragraph{Artifact pipeline.}
The prior-generation script
(\texttt{scripts/generate\_multimodel\_warmup\_priors.py}) uses the same
feature pipeline (Section~\ref{sec:cross_domain_transfer}): MiniLM embedding
$\to$ PCA-32 $\to$ bias append $=$ 33D context vectors.  The resulting
priors are stored in \texttt{priors\_warmup\_43model.joblib}; the split
definition is stored in \texttt{splits\_three\_way.json} for reproducibility.

\paragraph{Data provenance summary.}
Table~\ref{tab:data_provenance_extended} extends
Table~\ref{tab:data_provenance} to cover both the two-model and multi-model
experimental setups.

\begin{table}[h]
\centering
\caption{Extended data provenance covering two-model and multi-model experiments.}
\label{tab:data_provenance_extended}
\small
\begin{tabular}{llll}
\toprule
 & \textbf{PCA Training} & \textbf{$K{=}2$ Eval} & \textbf{$K{>}2$ Eval} \\
\midrule
\textbf{Source} & RouteLLM HuggingFace & LMSYS Chatbot Arena & LMSYS Chatbot Arena \\
\textbf{Models} & 2 (Mixtral, GPT-4-T) & 2 (same pair) & 43 models \\
\textbf{Size} & 80,000 prompts & 750 holdout & ${\sim}$533 online + 750 holdout \\
\textbf{Judge} & GPT-4 & GPT-4o & GPT-4o \\
\textbf{Split} & Independent dataset & Dev/holdout & Prior-train / online / holdout \\
\bottomrule
\end{tabular}
\end{table}

\subsubsection{Limitations}
\label{sec:transfer_limitations}

\paragraph{1. Two-model topology.}
The PCA is trained on Mixtral-vs-GPT-4-Turbo battles specifically.  The
variance it captures is aligned with \emph{this} capability boundary.  For a
different model pair, the routing-relevant directions in embedding space may
differ, and the PCA would need retraining.  The production library supports
JIT PCA retraining (\texttt{allow\_jit\_training=True}) to handle this.

\paragraph{2. Semantic near-duplicates.}
The 243 removed overlaps were detected via exact string matching.  Semantic
near-duplicates (paraphrases, translations, minor edits) may still exist
between the warmup and holdout sets.  Given the 80,000-to-750 size ratio and
independent provenance, any residual near-overlap is unlikely to meaningfully
inflate the routing signal, but we cannot rule it out entirely.

\paragraph{3. Stratified holdout.}
The holdout was constructed with a difficulty dimension derived from oracle
reward statistics, enriching it for prompts with clear model preferences.
This means the absolute magnitude of correlation statistics (Spearman $\rho$,
signal ratio) may be inflated relative to an unstratified deployment
distribution.  The \emph{existence} of the signal is robust; the
\emph{magnitude} in production depends on the user's prompt mix.

\paragraph{4. Embedding model dependency.}
The entire pipeline depends on \texttt{all-MiniLM-L6-v2} as the sentence
encoder.  A different encoder would produce different 384D representations,
invalidating the trained PCA.  The PCA and encoder must be treated as a
coupled pair.
